%% tables.tex
%% Copyright 2026 E. Nijenhuis  -- LPPL 1.3c, see LICENSE.txt.
%%
%% Edge-case tests for the `table' parameter type.  Each section exercises
%% one scenario where the previous string-replacement implementation either
%% failed or behaved surprisingly.  After the row-stack refactor, every
%% case below should render the same way it would if the user had hand-
%% written the tabular by inlining payload values.
\documentclass{article}
\usepackage{lua-placeholders}
\usepackage[autolanguage]{numprint}

\loadrecipe[\jobname]{tables-spec.yaml}
\loadpayload[\jobname]{tables-payload.yaml}

\setlength{\parindent}{0pt}
\nprounddigits{2}

\begin{document}

    \section*{Case 1: All columns referenced (baseline)}

    \newcommand\rowAB{\desc & \price \\}
    \begin{tabular}{l|l}
        \textbf{Description} & \textbf{Price} \\\hline
        \fortablerow{basic table}{rowAB}
    \end{tabular}

    \section*{Case 2: Subset of columns referenced}

    Omitting \textbackslash price\ in the row macro should drop the column
    silently --- the row still iterates and the description column still
    renders.  The unreferenced cell value is computed but never expanded.

    \renewcommand\rowAB{\desc \\}
    \begin{tabular}{l}
        \textbf{Description} \\\hline
        \fortablerow{basic table}{rowAB}
    \end{tabular}

    \section*{Case 3: Columns referenced in reverse order}

    With string replacement, swapping the order in the row macro could
    interfere when one column name is a prefix of another.  With the
    macro-binding approach, order is irrelevant.

    \renewcommand\rowAB{\price & \desc \\}
    \begin{tabular}{r|l}
        \textbf{Price} & \textbf{Description} \\\hline
        \fortablerow{basic table}{rowAB}
    \end{tabular}

    \section*{Case 4: Underscore column names with \texttt{\textbackslash ExplSyntaxOn}}

    The \texttt{underscore~table} declares three columns: \texttt{item}
    (a number), \texttt{item\_name} and \texttt{unit\_price}.  Two issues
    appear:

    \begin{itemize}
        \item \texttt{\textbackslash item} is a prefix of
              \texttt{\textbackslash item\_name}.  String-replacement
              substituted \texttt{\textbackslash item} first, mangling
              \texttt{\textbackslash item\_name}.  The macro-binding
              implementation defines each column as its own control
              sequence, so prefix collisions are impossible.
        \item Without \texttt{\textbackslash ExplSyntaxOn}, \TeX{} reads
              \texttt{\textbackslash item\_name} as
              \texttt{\textbackslash item} followed by the subscript
              token \texttt{\_} and the letters \texttt{name}, which
              fails.  Wrap the row macro in
              \texttt{\textbackslash ExplSyntaxOn~\dots
              \textbackslash ExplSyntaxOff} to flip \texttt{\_} to
              letter catcode at parse time so \texttt{\textbackslash
              item\_name} tokenises as one control sequence.
    \end{itemize}

    \ExplSyntaxOn
    \cs_new:Npn \rowU {\item & \item_name & \unit_price \\}
    \ExplSyntaxOff

    \begin{tabular}{r|l|r}
        \textbf{\#} & \textbf{Item} & \textbf{Unit price} \\\hline
        \fortablerow{underscore table}{rowU}
    \end{tabular}

    \section*{Case 5: Special characters in cell values}

    A cell may carry whatever \TeX{} can render --- but the YAML author
    still has to obey \TeX's lexical rules.  In particular, a literal
    \texttt{\%} must be written as \texttt{\textbackslash\%} in the
    payload, otherwise \TeX{} sees the unescaped \texttt{\%} and treats
    the rest of the line as a comment.  UTF-8 characters such as
    \texttt{€} pass through unchanged because Lua\TeX's input encoding is
    UTF-8 by default and the body font has the glyph.

    Old string-replacement was also vulnerable to a Lua-pattern hazard:
    \texttt{\%} in a cell value would be interpreted as a back-reference
    by \texttt{string.gsub}'s replacement argument.  With
    \texttt{token.set\_macro} the value is just bytes, so this class of
    breakage is gone.

    \newcommand\rowW{\label & \note \\}
    \begin{tabular}{l|l}
        \textbf{Label} & \textbf{Note} \\\hline
        \fortablerow{weird value table}{rowW}
    \end{tabular}

\end{document}
